\chapter{The Calibration Profile}\label{ch05:profile}

Every quantity in CAS-Ohm 1.0 is relative to a calibration profile. This
chapter develops that object field by field: what each field binds, why the
binding is mandatory rather than advisory, what failure the field's absence
would permit, and how the fields jointly support the two principles that
organize the whole candidate---\emph{profile relativity} (no quantity means
anything except under its profile, and nothing transfers across profiles
without a stated comparison rule) and \emph{anti-circularity} (an effective
resistance is fit on calibration data and earns its keep only by predicting
held-out data, developed fully in Chapter~\ref{ch06:response}).

The motivating observation is simple and easy to underestimate. A
directional activation projection is not a property of ``the model'' in the
abstract. It is a property of a specific weights file, read by specific
code, under a specific runtime and numeric precision, at a specific layer
and token position, projected onto a specific direction vector of specific
provenance, normalized against a full-scale value obtained by a specific
method from a specific corpus, with specific seeds. Change any element and
the number changes---sometimes negligibly, sometimes qualitatively, and
\emph{which} of those it will be is exactly what nobody knows in advance for
systems of this kind. Metrological practice long ago internalized the
consequence: a measurement result is inseparable from the procedure and
conditions that produced it, and a reported value without its conditions is
not a result but a rumor \cite{vim,gum}. The calibration profile is the
candidate's mechanization of that doctrine. It is a single, frozen,
content-addressed record binding every condition that could plausibly move a
number, so that ``the same measurement'' has an exact meaning: same profile
digest, nothing less.

The chapter's claim discipline matches the rest of the monograph. The
profile is an implemented data contract with validation rules exercised by
bounded software tests; no live profile for any production model has been
admitted, no field's adequacy has been empirically validated, and the whole
apparatus is \textsc{ring 4 / not licensed / held}. What is defended here is
the \emph{design argument}: why each field exists and what its omission
would break.

\section{The Profile as the Unit of Meaning}\label{ch05:unit}

\begin{definition}[Calibration profile]\label{ch05:def:profile}
A calibration profile is an immutable record binding: a profile identifier;
the exact model identity (model, configuration, tokenizer, weights); the
code revision and dependency lock; the runtime identity and dtype; temporal
validity (creation, validity start, expiry) and drift state with its
checking evidence; the registered test points; the direction vectors with
identity, norm, sign, and provenance; the corpus manifest and its rights
record; the calibration and validation split references; the random seed and
randomization method; the alpha grid with its endpoints; the full-scale
($I_{\mathrm{FS}}$) method, its per-test-point full-scale values, and, where
applicable, its percentile parameter; the fit method and its receipt; the
uncertainty method; the response-region method; the tolerance profile; the
control references and per-test-point control bindings; the source
references; and a review state. The profile's canonical digest is the
SHA-256 of its canonical encoding, and every measurement made under the
profile binds that digest.
\end{definition}

\begin{invariant}[Profile relativity]\label{ch05:inv:relativity}
No CAS-Ohm quantity is interpretable except relative to the calibration
profile whose digest it binds. No quantity, threshold, region boundary,
full-scale value, fitted resistance, or classification transfers to another
profile---including a profile differing in a single field---without an
explicitly stated and separately evidenced comparison rule.
\end{invariant}

The invariant sounds austere; it is in fact merely honest. Two profiles that
differ in, say, dtype may produce projections that differ in the third
decimal or in the first, and absent a comparison experiment no one knows
which. Declaring the numbers comparable by default would be an empirical
claim made without evidence. The candidate's default is therefore
incomparability, with comparability available as a first-class,
evidence-bearing assertion when someone does the work to establish it. This
inverts the usual failure mode of measurement systems, in which comparability
is assumed until an incident proves otherwise; here incomparability is
assumed until a comparison rule is admitted. Chapter~\ref{ch04:audit} showed
why the inversion is forced: the sources themselves demonstrate
(Conflict~\ref{ch04:conf:ifs}) that even the framework's own full-scale
denominators were never one quantity.

A second consequence of Definition~\ref{ch05:def:profile} deserves emphasis
before the field-by-field walk: the profile is \emph{immutable}. There is no
``update a profile'' operation. A changed condition---new weights, a revised
corpus, a different tolerance---produces a new profile with a new digest,
and history accumulates as a set of distinct profiles rather than as
mutations of one. Immutability is what makes the digest meaningful: a digest
that could survive edits would identify nothing.

\section{Identity of the Measured System}\label{ch05:model}

The profile's model identity binds five values: a model identifier, a
revision string, and three content digests---weights, configuration, and
tokenizer. The rationale for the identifier and revision is ordinary
bookkeeping; the rationale for the three digests is that names and revisions
are assertions while digests are checkable facts. Model repositories reuse
names across incompatible checkpoints; configurations change activation
functions and attention variants without renaming the model; tokenizers
change vocabulary and merge rules in ways that silently re-index every token
position a test point names. Each of the three artifacts can vary
independently of the other two, and each variation invalidates measurements
in a different way: changed weights change every activation; a changed
configuration can change the meaning of ``layer 12''; a changed tokenizer
changes which token sits at position 7 of a probe prompt, which is enough to
re-aim every token-position-bound test point at different content.

The profile also requires that its model artifact references \emph{resolve}:
the set of digests among the artifact references must equal the set of
digests named in the model identity. This closes a subtle gap between
naming and having. A profile that named a weights digest without a
resolvable artifact reference would be unfalsifiable---nobody could re-derive
its measurements, because nobody could obtain its weights. The resolution
rule makes the profile self-supporting: everything it names, it references,
and the evidence registry can demand the bytes.

\section{Code Revision, Dependency Lock, Runtime, and Dtype}\label{ch05:code}

Measurements are produced by code, and code has the same identity problems
as models. The profile binds a code reference (the exact revision of the
measuring implementation) and a dependency lock reference (the closed set of
third-party versions beneath it). The dependency lock matters because the
measuring code's behavior is not a function of its own source alone:
numerical libraries change reduction orders, kernel selections, and default
tolerances between versions, and any of these can move low-order digits of
an activation projection. A measurement that binds its code revision but not
its dependency closure has fixed the recipe while leaving the ingredients
loose. The lock reference freezes the closure; re-execution under a
different closure is, by
Invariant~\ref{ch05:inv:relativity}, a different profile.

The runtime identity and dtype fields record the execution substrate and the
numeric precision of the computation. Dtype is the more consequential of
the two and the easier to lose track of, since frameworks freely and
sometimes implicitly cast between precisions. Half-precision, bfloat16, and
single-precision evaluations of the same network on the same input are three
different computations with three different rounding profiles; whether their
differences matter for a given direction projection is an empirical question
that cannot even be \emph{asked} unless the precision is recorded. The
profile makes the recording mandatory, so that the question is always
askable and never silently begged.

\section{Instruments}\label{ch05:instruments}

The profile binds a non-empty tuple of instrument identities---each carrying
an instrument identifier, a kind, a version, and an implementation content
digest---together with implementation references that must resolve exactly
the set of implementation digests the instruments name. An instrument, in
the candidate's vocabulary, is the identified software chain that produces
an instrumented measurement: the hooking and projection code that reads an
activation, applies the direction, and emits the raw value. The term is
borrowed deliberately from laboratory practice, and so is the discipline: a
laboratory result is attributed to a specific instrument with a calibration
binding, not to ``the lab'' \cite{vim}.

Two rationales drive the field. First, the measurement-kind rules of the
ontology make \texttt{INSTRUMENTED} the most privileged kind---the only kind
that can anchor structural research-packet completeness---and privilege must
be checkable: an instrumented measurement must name a registered instrument,
and the profile must hold the bytes of that instrument's implementation, or
the claim ``this number came from an identified instrument chain'' is
unverifiable decoration. Second, instruments version independently of the
main code revision. The measuring harness may be stable across many profile
revisions while a projection kernel changes, or vice versa; giving
instruments their own identities and digests means the question ``did the
instrument change between these two profiles?'' has a mechanical answer,
field by field, rather than requiring a code archaeology expedition.

The uniqueness rule (instrument identifiers unique within a profile) and
the resolution rule (implementation references exactly matching declared
digests) mirror the model-artifact discipline of
Section~\ref{ch05:model}, and for the same reason: naming without having is
the failure mode, and exact set equality between names and holdings is the
check.

\section{Temporal Validity, Drift, and Expiry}\label{ch05:drift}

A profile carries three timestamps---created at, valid from, expires
at---ordered by the validation rule (creation not after validity start,
validity start strictly before expiry), plus a drift state and a drift-check
evidence reference. Together they implement a position the candidate takes
deliberately: \emph{calibration decays}. Nothing about a frozen weights file
decays, of course; what decays is the environment of assumptions around
it---upstream serving stacks change, corpora age relative to usage,
directions estimated from one data distribution meet another. A calibration
regime without expiry drifts into exactly the ``stale profile'' condition
that Chapter~\ref{ch04:audit}'s residual discipline identified as a standing
alternative explanation for anomalies. The profile therefore expires by
construction, and expiry is checked mechanically: the currency assertion
fails outside the validity window.

The drift state is a four-valued enumeration---\texttt{CURRENT},
\texttt{STALE}, \texttt{FAILED}, \texttt{REVOKED}---bound to evidence by the
drift-check reference. \texttt{CURRENT} asserts a passing recent check;
\texttt{STALE} that the check is overdue; \texttt{FAILED} that a check ran
and did not pass; \texttt{REVOKED} that the profile was withdrawn by
decision. The distinction between \texttt{STALE} and \texttt{FAILED}
preserves an important asymmetry of ignorance: an overdue check means
\emph{we do not know} whether calibration holds, while a failed check means
\emph{we know it does not}. Both block operation (only \texttt{CURRENT}
passes the currency assertion), but they license different remediations and
different retrospective doubts about recent results.

Operationally, the profile's currency assertion requires the conjunction:
review state \texttt{REVIEWED}, drift state \texttt{CURRENT}, and the
replay time inside the validity window. All admitted operations evaluate it
against a caller-supplied replay time---which, as
Chapter~\ref{ch04:audit}'s eligibility discussion notes, is a historical
replay coordinate, never an authenticated assertion about the actual
present.

\section{Test Points}\label{ch05:testpoints}

A test point is the profile's answer to ``measured \emph{where}?'' Each
registered test point binds: the model identity (which must equal the
profile's model---a test point cannot smuggle in a different system); the
layer; the token positions (a non-empty tuple of distinct non-negative
indices); the direction vector's identifier, content digest, provenance,
norm, and sign; and the hook type (where in the layer's computation the
activation is read). Test points must be unique within a profile, and every
per-test-point structure elsewhere in the profile---full-scale values,
control bindings---must cover the registered test points exactly.

The rationale is locality. CAS-Ohm's response model is explicitly a
\emph{local} hypothesis: a dose-response relation at a particular place in a
particular computation, not a global property of a model. All of the audit's
warnings about transfer (full-scales that ``do not transfer,'' flow that
``remains local'') are warnings that locality is real and unforgiving. The
test point is locality made into an identity: two measurements are at the
same place if and only if their test-point digests are equal, which requires
agreement on layer, positions, direction, sign, and hook type
simultaneously. Anything less---``layer 12, roughly the same place''---is
not an identity, and the type system does not accept it as one.

The hook type deserves a word, since it is the field most often left
implicit in informal practice. A transformer layer is not a point; it is a
pipeline of sublayer computations, and ``the activation at layer 12'' could
mean the residual stream before the block, after attention, after the MLP,
or after the final layer norm---numerically different vectors with different
relationships to the direction vector. Recording the hook type turns an
ambiguous phrase into an exact coordinate.

\section{Direction Vectors: Identity, Norm, Sign, and Provenance}\label{ch05:directions}

The direction vector is the most epistemically delicate object in the
profile. Every response magnitude in the framework is a projection onto some
direction; the direction, in effect, \emph{is} the instrument's aperture.
The profile binds four aspects.

\emph{Identity.} Each direction carries an identifier and a content digest,
and the profile's direction references must resolve exactly the digests
named by its test points---the same resolution discipline as for model
artifacts, for the same reason: a direction that cannot be re-obtained
byte-for-byte cannot anchor a reproducible measurement.

\emph{Norm.} The candidate fixes direction norm to exactly one (validated as
\texttt{Decimal("1")}, an exact equality, not a tolerance). Projections
scale linearly with the direction's norm, so unnormalized directions would
fold an arbitrary multiplicative constant into every response value and
every fitted resistance. Fixing the norm removes a gauge freedom that would
otherwise make numbers incomparable even within one profile.

\emph{Sign.} A direction and its negation define the same line; a
projection onto one is the negative of the projection onto the other. Sign
conventions are thus pure convention---but an \emph{unrecorded} convention
is a latent sign error, and a sign error in a response model flips the
apparent semantic outcome (\texttt{INTENDED} versus \texttt{OPPOSITE})
wholesale. The profile records the sign explicitly so that the convention is
an auditable fact rather than tribal knowledge.

\emph{Provenance.} The candidate enumerates three direction provenance
kinds, and the enumeration is closed:

\begin{description}
\item[\texttt{PROBE\_DERIVED}.] The direction is the weight vector of a
probe---typically a linear classifier---trained to separate labeled
activation examples. Its epistemic character: the direction is an
\emph{estimate from data}, inheriting the probe corpus, its labels, its
regularization, and its sampling noise. Probe directions answer ``along
which axis do these labeled classes separate?''---which is related to, but
not the same as, ``along which axis does the model \emph{use} this
concept?''
\item[\texttt{JACOBIAN\_DERIVED}.] The direction is obtained from
differentiation---a gradient or Jacobian row of some downstream quantity
with respect to the activation. Its character: the direction is a
\emph{local analytic object}, exact for the computation at the linearization
point but valid only there; it answers ``what perturbation most moves this
output, here?''
\item[\texttt{DIFF\_DERIVED}.] The direction is a difference of activation
statistics between two conditions (for example, means under contrasting
prompts). Its character: a \emph{contrast estimate}, simple and robust, but
credited with meaning only to the extent the two conditions differ in one
thing.
\end{description}

Three kinds, three different failure surfaces: probe directions fail with
their training distribution, Jacobian directions fail away from their
linearization point, and diff directions fail when the contrast was
confounded. The profile does not rank them---each is legitimate---but it
refuses to let them blur, for the by-now familiar reason: a number whose
production method is unrecorded cannot be doubted \emph{specifically}, and
undoubtable numbers are the enemy. The provenance field also feeds the
registered reproducibility obligations (cross-split direction
reproducibility appears among the closure experiments), which are stated
per-provenance because the reproduction procedure differs by kind.

A corollary worth stating explicitly: directions of different provenance are
not interchangeable even when they point in nearly the same place. A probe
direction and a diff direction for ``the same concept'' at the same test
point are estimates of related but distinct estimands, with distinct error
structures; a resistance fitted against projections on one is not a
resistance for the other, and the profile machinery enforces this
automatically because the direction digest participates in the test-point
identity, which participates in everything downstream. If two provenances'
directions are claimed equivalent for some purpose, that claim is a
comparison rule in the sense of Invariant~\ref{ch05:inv:relativity}---an
assertable, evidence-bearing statement---not a default. The candidate's
posture is uniform: geometric proximity is measurable, semantic
equivalence is a claim, and claims need evidence.

\section{Corpus Manifest and Rights}\label{ch05:corpus}

Two references bind the data environment: the corpus manifest and the corpus
rights record. The manifest is a content-addressed inventory of the corpus
from which calibration stimuli, probe data, and (under the percentile
method) full-scale statistics are drawn. Content addressing does for the
corpus what digests do for weights: ``the corpus'' stops being a mutable
directory and becomes an exact set of bytes, so that corpus-derived
statistics---above all the 95th-percentile full-scale of
Conflict~\ref{ch04:conf:ifs}---are anchored to something that cannot
silently gain or lose documents.

The rights reference is the field most likely to strike a technical reader
as out of place, and its inclusion is deliberate. A calibration built on
data whose use rights are unestablished is a calibration that may not be
lawfully re-executed, shared with a reviewer, or retained---which makes
every downstream measurement irreproducible in the practical sense that
matters, whatever its technical merits. Binding a rights record into the
profile digest makes rights status part of measurement identity: a profile
cannot be complete without it, and a rights change (a revoked license, a
newly discovered encumbrance) is visible as exactly what it is, a change to
the conditions under which the numbers exist.

\section{Splits, Seeds, and Randomization}\label{ch05:splits}

The profile binds a calibration split reference and a validation split
reference, and its validation rules require the two digests to differ. This
pair of fields is the structural footprint of the anti-circularity rule: the
data on which the effective resistance is fit (calibration split) and the
data on which its predictions are scored (validation split) are distinct
artifacts, distinguished at the type level, before any statistics happen.
Chapter~\ref{ch06:response} develops the statistical argument; here we note
only the design point that the split separation is a \emph{profile} fact,
not a habit of the analyst. An analyst cannot forget the held-out discipline
without constructing a profile that fails validation.

The random seed (a non-negative integer) and the randomization method (an
identifier for how randomness was consumed---assignment of items to splits,
ordering of sweeps, any stochastic element of the fit) close the remaining
degrees of freedom. Randomization without a recorded seed is
irreproducibility by construction; a seed without a recorded method is a
seed for an unspecified generator consumed in an unspecified order, which
reproduces nothing. Together they make the stochastic parts of calibration
replayable, and---as important---they make \emph{selective} reruns visible:
a report whose seed differs from its profile's seed is advertising that
something was rerun.

\section{The Alpha Grid}\label{ch05:alphagrid}

The alpha grid is the ordered tuple of intervention doses at which the
dose-response relation is probed. The validation rules are specific: at
least three doses; exact decimal values; strictly increasing and duplicate
free; the declared minimum and maximum equal to the grid's endpoints; and
zero included.

Each clause has a job. \emph{Three doses minimum}: two points determine a
line exactly, leaving zero residual degrees of freedom---no way to detect
curvature, hence no way to distinguish a genuinely linear response from any
monotone response whatsoever; three is the least count at which
``linearity'' is even falsifiable, and the region classifier's own minimum
(interior points with neighbors on both sides) matches it. \emph{Exactness
and strict ordering}: the grid is part of measurement identity, and a grid
with representation-dependent values or ambiguous ordering would leak
platform noise into the profile digest. \emph{Endpoints declared}: the
declared range is the region over which the profile claims to have probed;
predictions outside it are extrapolations, and the admissibility machinery
can only know that if the range is a recorded fact. \emph{Zero included}:
the zero dose is the null intervention---the control column built into the
grid itself. Without it, the fit has no anchor at the origin and the
zero-intervention control comparisons required by the closure experiments
have no in-grid reference.

The grid also encodes a modeling humility: it is a \emph{grid}, not a
distribution. The candidate does not assume responses at unprobed doses; the
response region machinery classifies behavior at and between probed doses,
and the tolerance profile (Section~\ref{ch05:tolerances}) governs how much
interpolation credence the fit may claim.

\section{The Full-Scale Method}\label{ch05:ifs}

Conflict~\ref{ch04:conf:ifs} established that the sources supply three
incompatible procedures for the full-scale reference $I_{\mathrm{FS}}$. The
profile's response is the \texttt{IFSMethod} enumeration, with the audit's
disposition implemented literally: one method per profile, declared up
front, inside the digest.

\begin{description}
\item[\texttt{RAIL\_AND\_FLOOR\_PER\_PATH}.] The Appendix~I procedure:
$I_{\mathrm{FS}}$ measured once per path at drive rail and constraint
floor. An extreme-operating-point probe; local by construction, with the
source's own warning that it does not transfer.
\item[\texttt{CORPUS\_PERCENTILE\_PER\_TEST\_POINT}.] The addendum
procedure: the per-test-point 95th percentile of absolute directional
projection over a fixed hashed corpus. A statistic, requiring the
percentile parameter; the profile's validation rules enforce that the
percentile is present exactly when this method is declared and lies
strictly between 0 and 100.
\item[\texttt{DECLARED\_DESIGN\_FULL\_SCALE}.] The Guide procedure: a
board-selected design value. A declaration rather than a measurement, typed
accordingly (its natural measurement kind is a design value, not an
instrumented one).
\end{description}

\begin{ruleenv}[Non-interchangeability of full-scale methods]\label{ch05:rule:ifs}
The three \texttt{IFSMethod} values name different calibration procedures
producing different quantities. No value obtained under one method may be
substituted for, averaged with, or compared against a value obtained under
another except through an explicitly admitted comparison rule. Because the
method is bound into the profile digest, any cross-method comparison is a
cross-profile comparison and is blocked by default under
Invariant~\ref{ch05:inv:relativity}.
\end{ruleenv}

Whatever the method, the profile must carry a full-scale \emph{value record}
for every registered test point, exactly one each. The value record is
itself typed evidence: a strictly positive value in activation units, an
uncertainty whose interval is positive and contains the value, a sample
size, and non-empty source references. Note the asymmetry with the method
enumeration: the method says how the denominator was obtained; the value
record proves the denominator exists, with uncertainty, for each place it
will be used. Normalization (Chapter~\ref{ch08:quantization} where present)
divides by these values and propagates their uncertainty into every Q10
score; a full-scale without uncertainty would launder the denominator's
error out of every displayed value.

\section{Fit Method, Uncertainty Method, and the Fit Receipt}\label{ch05:fit}

Three fields govern how the effective resistance is estimated: the fit
method identifier, the uncertainty method identifier, and the fit receipt
reference. The method identifiers are deliberately identifiers rather than
enumerations: the candidate does not legislate a unique estimator (least
squares on the active region, robust variants, and others are all
conceivable), but it insists the choice be named, because estimator choice
moves results and an unnamed estimator is an unfalsifiable one. The
uncertainty method is bound separately because uncertainty computation is a
procedure in its own right (analytic propagation, bootstrap, and interval
methods differ), and the audit's discipline requires every emitted
uncertainty to say how it was obtained \cite{gum}.

The fit receipt reference binds the artifact recording the actual fit run:
inputs consumed, estimates produced, diagnostics observed. Its presence in
the profile closes the loop between declaration and execution---the profile
does not merely promise a fit method, it references the receipt of the fit
that was actually performed. In the Kennedy--O'Hagan framing of computer
model calibration, the receipt is where the calibration data, the parameter
posterior (or point estimate with its uncertainty), and the discrepancy
evidence live \cite{kennedyohagan}; Chapter~\ref{ch06:response} returns to
what the receipt must show before a resistance earns the
\texttt{CALIBRATED\_HELD\_OUT} role.

\section{Tolerances}\label{ch05:tolerances}

The tolerance profile is a five-field record of exact decimals: an absolute
tolerance, a relative tolerance, a residual floor, a minimum linearity
coefficient of determination, and a maximum off-target fraction. Validation
requires the absolute tolerance and residual floor strictly positive and the
relative tolerance, linearity minimum, and off-target maximum in the unit
interval.

The rationale is that every downstream judgment of the form ``close
enough,'' ``linear enough,'' or ``clean enough'' must be a comparison
against a declared number rather than an analyst's mood. The absolute and
relative tolerances govern agreement checks between predicted and observed
response. The residual floor sets the magnitude below which a residual is
treated as indistinguishable from zero---the operational meaning of
``zero-compatible'' in the residual-censoring rule of
Conflict~\ref{ch04:conf:residual}. The linearity minimum is the floor an
active-region fit must clear before any operation gated on ``adequate
linearity'' (notably the network compositions of
Conflict~\ref{ch04:conf:networks}) may proceed. The off-target maximum
bounds the fraction of response appearing on axes other than the target
direction before the measurement is disqualified as off-axis-contaminated.

Putting the five numbers in the profile---rather than in code defaults or
per-analysis choices---has two effects. It makes tolerance-shopping visible:
loosening a tolerance to admit a marginal result requires minting a new
profile with a new digest, an auditable event. And it makes strictness
comparable: two profiles' tolerance records can be diffed field by field,
so ``profile B is laxer than profile A'' is a checkable statement.

\section{Response-Region Rules}\label{ch05:regions}

The response-region method identifier names the procedure by which sweep
observations are classified into the dose-response regions
(\texttt{CUTOFF}, \texttt{ACTIVE}, \texttt{SATURATION}, \texttt{BREAKDOWN},
\texttt{NON\_OHMIC}, \texttt{UNSTABLE}, with \texttt{UNASSESSED} as the
honest default). The field exists in the profile for the same reason as the
fit method: region assignment is a procedure with alternatives, region
labels gate consequential operations (predictions are admissible only in the
active region; network compositions require active-region components;
saturated-stage residuals are censored), and a gate whose criterion is not
frozen is not a gate. Binding the method into the digest means a
reclassification standard cannot drift under a fixed profile: the same
sweep, replayed, classifies the same way.

\section{Controls and Control Bindings}\label{ch05:controls}

The profile binds a set of control references and a set of per-test-point
control bindings, with validation rules enforcing exact correspondence: the
bound control artifacts must exactly equal the declared control references,
every binding must name a registered test point, and the bindings must
cover every test point. Controls here are the comparison conditions the
closure experiments require---zero-intervention runs, random same-norm
directions, wrong-layer and wrong-token probes, label shuffles, norm-matched
perturbations---represented as artifacts with relevance methods binding them
to the specific test points they control for.

The design point is that controls are part of the \emph{calibration},
not an optional garnish on the analysis. A measurement under the profile
carries the profile's control set; the validation machinery rejects
non-design measurements whose control references do not equal the profile's.
This forecloses the quiet failure in which controls exist somewhere in a
repository but the measurement that mattered was taken without them. The
coverage rule (every test point controlled) similarly forecloses the partial
version, where the well-behaved test points get controls and the interesting
one does not.

\section{Review State}\label{ch05:review}

The profile carries a review state: \texttt{UNASSESSED}, \texttt{IN\_REVIEW},
or \texttt{REVIEWED}; the implementation additionally provides
\texttt{FLAG\_FOR\_HUMAN\_REVIEW} for records requiring escalation. Only
\texttt{REVIEWED} passes the currency assertion---an unreviewed profile can
be constructed, stored, and studied, but no operation that requires a
current profile will run under it.

The field encodes a separation the candidate is careful about everywhere:
\emph{structural validity is not endorsement}. A profile that passes every
validation rule is well formed; whether its choices are \emph{sensible}---is
this corpus appropriate, are these tolerances defensible, is this direction
provenance suited to this question---is a judgment no type system can make.
The review state is the socket where human judgment plugs into the
machinery, and its default is the honest \texttt{UNASSESSED}. Note what the
field does not claim: \texttt{REVIEWED} records that review occurred within
whatever process governs the record; it is not an independent-authority
attestation, which belongs to the separate authority machinery of
Chapter~\ref{ch11:eligibility} where present, and which the candidate
implements with constant-false authenticated eligibility.

\section{The Canonical Digest}\label{ch05:digest}

Every field above participates in the profile's canonical payload, which is
serialized by the candidate's canonical encoding (exact-type tagged,
cycle-aware, bounded) and hashed to produce the profile's SHA-256 digest.
The digest is the profile's identity for every purpose: measurements bind
it, operation contexts snapshot it and re-verify it on every subsequent
operation (the context fails closed if the profile, its authority manifest,
its attestations, or the replay time change after admission), and
cross-profile comparisons are detected by digest inequality.

Three design properties make the digest trustworthy. \emph{Totality}: the
payload includes every semantically relevant field---including the status
constants (admission ring, license status, disclosure status) and the
schema label---so no meaningful variation escapes the hash. \emph{Canonical
form}: field ordering and collection ordering are normalized (references
sorted by digest, bindings by their own hashes), so equality of content
implies equality of bytes implies equality of digest, with no
representation-dependent false negatives. \emph{Exactness}: numeric fields
are exact decimals, not binary floats, so the digest cannot be perturbed by
platform-dependent rounding in the serialization path. The result is that
``same profile'' is decidable by string comparison, which is the property
every downstream fail-closed check relies on.

\section{The Operation Context: Admission and Revalidation}\label{ch05:context}

The profile does not act alone; it acts inside an \emph{operation context},
the runtime object through which every measurement validation, normalization,
and structural operation flows. A context binds the profile, the evidence
registry holding the referenced bytes, the authority configuration with its
trust mode, any capability attestations, the caller-supplied replay time,
and optionally a clock-receipt record. At admission, the context performs
the full battery: it revalidates the evidence registry, resolves the
authority manifest's bytes against its digest, asserts the profile current
at the replay time (review state \texttt{REVIEWED}, drift
\texttt{CURRENT}, replay time inside the validity window), and demands that
every artifact the profile references---model artifacts, instrument
implementations, directions, splits, corpus, rights, controls, sources,
full-scale sources, the fit receipt---be present and resolvable. Then it
snapshots: the profile digest, the manifest digest, the attestation digests,
the replay time, and the clock-receipt digest are recorded, and every
subsequent operation re-verifies the snapshot before doing anything else.
Any drift---a profile object swapped, an attestation added, a replay time
edited---fails the operation closed.

The design rationale is time-of-check versus time-of-use. A validation
performed once at construction protects nothing if the validated objects can
change before use; the snapshot-and-revalidate pattern makes the context's
guarantees continuous rather than momentary. For the audit posture of this
monograph, the significant property is that the context is where profile
discipline becomes \emph{unavoidable}: there is no API path to a
normalization or a structural operation that does not pass through a
context, hence through the currency assertion, hence through the profile's
review, drift, and validity gates. The context also fixes the epistemic
status of time: the replay time is caller-supplied history, useful for
replaying past states, and the context's authenticated-currency answer is
constant false---no admitted external time authority exists in the
candidate, and the machinery says so rather than pretending.

\section{What Each Field Defends Against}\label{ch05:threats}

A compact way to review the chapter is by threat: for each family of
fields, the failure it renders unrepresentable or fail-closed.
Table~\ref{ch05:tab:threats} gives the mapping.

\begin{table}[htbp]
\centering\small
\begin{tabular}{@{}p{0.42\linewidth}p{0.5\linewidth}@{}}
\toprule
Field family & Failure defended against \\
\midrule
Model identity + artifact resolution & silent checkpoint/config/tokenizer
substitution; unobtainable weights behind a name \\
Code ref + dependency lock & behavior drift from unpinned libraries;
``same code'' with different numerics \\
Runtime + dtype & precision-dependent results with no record of precision \\
Instruments + implementation refs & privileged \texttt{INSTRUMENTED} claims
with unverifiable production chains \\
Validity window + drift state + check ref & stale calibration operating as
current; unknown-vs-known-bad conflation \\
Test points + hook type & ``roughly the same place'' passing as identity \\
Direction identity/norm/sign/provenance & gauge freedom in magnitudes; sign
flips in outcomes; method-blind pooling of estimates \\
Corpus manifest + rights & mutable corpora under fixed statistics;
legally unre-executable calibrations \\
Distinct splits + seed + method & circular fitting; unreplayable randomness;
silent selective reruns \\
Alpha grid rules & unfalsifiable linearity; undeclared extrapolation;
missing null dose \\
\texttt{IFSMethod} + full-scale records + percentile rule &
denominator-method conflation (Conflict~\ref{ch04:conf:ifs});
uncertainty-free denominators \\
Fit/uncertainty/region method ids + receipt & unnamed estimators; bare point
estimates; drifting region criteria; promised-but-unperformed fits \\
Tolerance profile & analyst-relative ``close enough''; invisible
tolerance-shopping \\
Controls + bindings + coverage & control-free measurements; selectively
controlled test points \\
Review state & structural validity mistaken for endorsement \\
Canonical digest & identity by name rather than content; representation
noise in identity \\
\bottomrule
\end{tabular}
\caption{Profile fields and the failures they foreclose.}
\label{ch05:tab:threats}
\end{table}

\begin{threat}\label{ch05:threat:residue}
The table's coverage claim should itself be read with the monograph's
discipline: these are the failures the fields are \emph{designed} to
foreclose, at the data-contract level, as exercised by bounded software
tests. Whether the fields suffice against adversarial construction of
profiles, against collusion between a profile author and a reviewer, or
against failure modes not yet enumerated is unassessed; the authenticated
authority layer that would address the adversarial cases is exactly the
part the candidate implements as constant-false eligibility
(Chapter~\ref{ch11:eligibility} where present).
\end{threat}

\section{An Illustrative Construction}\label{ch05:construction}

\begin{construction}[Hypothetical profile walk-through; all values
illustrative]\label{ch05:constr:example}
The following is a fabricated example to make the field interactions
concrete. Every value is invented for exposition; nothing here describes a
real model, corpus, or measurement.

An analyst wishes to calibrate a response model at one test point of a
hypothetical checkpoint. She constructs: a model identity naming the
checkpoint with its three digests, and three artifact references resolving
them; a code reference at her harness revision and a dependency lock from
her environment export; runtime identifier and dtype
(say, \texttt{"runtime-x"} and \texttt{"float32"}); a validity window of
ninety days from a drift check whose passing record she references, drift
state \texttt{CURRENT}; one instrument identity for her projection hook
with its implementation digest resolved; one test point at layer
\texttt{"blocks.12.hook\_resid\_post"}, token positions $(5, 9)$, with a
\texttt{DIFF\_DERIVED} direction of unit norm, positive sign, and resolved
digest; corpus and rights references for her fixed prompt set; two split
references partitioning it (digests distinct, as required); seed $17$ with
a named shuffling method; an alpha grid
$(-1.0, -0.5, 0, 0.5, 1.0)$ with endpoints declared and zero present;
\texttt{IFSMethod.CORPUS\_PERCENTILE\_PER\_TEST\_POINT} with percentile
$95$ and one full-scale record for her test point (value, uncertainty
interval, sample size, sources); fit method
\texttt{"ols-active-region-v1"} with a receipt reference; uncertainty
method \texttt{"bootstrap-1000-v1"}; region method
\texttt{"piecewise-slope-v1"}; tolerances (absolute $0.05$, relative
$0.10$, residual floor $0.02$, linearity minimum $0.95$, off-target
maximum $0.20$); control references and bindings covering the test point
(zero-dose, random-direction, wrong-layer artifacts); source references to
her design notes; review state \texttt{UNASSESSED}.

The profile validates structurally---every rule of this chapter
passes---and receives its digest. Note what she can and cannot now do. She
cannot run a normalization: the currency assertion fails on
\texttt{UNASSESSED}. After review, with state \texttt{REVIEWED} (a new
profile, hence a new digest---immutability), operations under an admitted
context become possible. If she later regenerates her direction vector,
even ``equivalently,'' its digest changes, so the test point changes, so
the profile changes; her earlier full-scale record no longer covers the
new test point, and validation forces her to remeasure it. The contract has
done its job: an easy-to-overlook invalidation (new direction, old
denominator) was surfaced as a hard error at construction time rather than
as a subtly wrong number at analysis time.
\end{construction}

\section{Profile Relativity in Practice}\label{ch05:relativity}

It is worth spelling out what Invariant~\ref{ch05:inv:relativity} concretely
forbids and what it permits, because the invariant is the profile's entire
point.

Forbidden by default: comparing Q10 scores across profiles (different
denominators, per Rule~\ref{ch05:rule:ifs}); pooling sweep observations
across profiles into one fit; carrying a fitted $r_{\mathrm{eff}}$ from one
model revision to the next (``the model barely changed'' is an empirical
claim requiring evidence, not a default); quoting a region classification
under a profile whose drift state is not current; and every flavor of
cross-model league table---``model A has lower resistance than model
B''---absent an admitted comparison rule. The last of these is the
politically hardest to resist, since ranking is what audiences want; it is
also the least defensible, compounding every incomparability at once
(different weights, tokenizers, geometries, direction provenances, and
full-scale denominators).

Permitted: anything within one digest, subject to the other gates; and
cross-profile statements made \emph{through} an explicitly admitted
comparison rule---a registered procedure with its own evidence establishing
that a specific quantity is comparable across a specific pair of profiles
under specific conditions. The registered checkpoint-replication hypothesis
is an instance: replicate on a held-out checkpoint \emph{under a
prespecified comparison rule}. The invariant does not make cross-profile
science impossible; it makes it a deliberate, evidenced act instead of an
ambient assumption.

\begin{remark}\label{ch05:rem:units-analogy}
Profile relativity generalizes the familiar rule that quantities in
different units may not be added without conversion \cite{buckingham}. A
profile is, in effect, a very large compound unit: the ``unit'' in which a
CAS-Ohm number is expressed includes the weights digest, the direction, the
denominator method, the dtype, and the rest. Two numbers in different
compound units need a conversion---the comparison rule---and the burden of
producing it lies with whoever asserts comparability. The type system merely
refuses to pretend the conversion is the identity.
\end{remark}

\section{The Profile Against the Audit}\label{ch05:vsaudit}

The profile can also be read as the point-by-point answer to
Chapter~\ref{ch04:audit}. Each preserved conflict deposits at least one
field.

Conflict~\ref{ch04:conf:scale} (scale and saturation) appears as the
insistence on exact decimal values everywhere and on typed quantities bound
to unit kinds; the profile contains no field where a number floats free of
its scale. Conflict~\ref{ch04:conf:ifs} (full-scale profile) deposits the
\texttt{IFSMethod} enumeration, the percentile-presence rule, and the
per-test-point full-scale records with mandatory uncertainty---the audit's
``separate named profiles'' realized as literally separate digests.
Conflict~\ref{ch04:conf:bridge} (dimensional bridge) deposits the entire
fit apparatus: the split references, fit method, fit receipt, and
uncertainty method exist so that an alpha-per-activation quantity can come
into existence by calibration rather than by relabeling a score.
Conflict~\ref{ch04:conf:gain} (gain definition) deposits the method-identity
discipline---dimensionless quantities carry method identifiers because unit
checking cannot distinguish them. Conflict~\ref{ch04:conf:asom} (ASOM axes)
deposits the response-region method field, which fixes the region axis's
procedure while leaving the semantic axis to independent assessment.
Conflict~\ref{ch04:conf:networks} (network rules) deposits the tolerance
profile's linearity minimum and the region method, which are the
preconditions the capability gates consult.
Conflict~\ref{ch04:conf:residual} (residual identification) deposits the
residual floor---the operational meaning of ``zero-compatible''---and the
control bindings that make shuffled and matched controls a structural
requirement rather than an aspiration. Conflict~\ref{ch04:conf:q4} (Q4
boundaries) deposits the exactness discipline itself: band assignment is
computable from exact pre-rounding values only because the profile's
numeric fields are exact from end to end.

This mapping is the concrete sense in which the audit ``compiles'' into the
candidate: the conflicts of Chapter~\ref{ch04:audit} are not commemorated in
a document appendix but disassembled into validation rules that run on
every profile ever constructed.

\section{Anatomy of a Comparison Rule}\label{ch05:comparison}

Invariant~\ref{ch05:inv:relativity} repeatedly defers to ``a stated
comparison rule.'' Since the deferral carries so much weight, the candidate
owes a statement of what such a rule must contain, even though none has yet
been admitted. A comparison rule, to be admissible, must at minimum
declare: (i) the two profile digests it relates, exactly---not families,
not patterns; (ii) the specific quantity kinds it renders comparable, with
all others remaining incomparable; (iii) the transformation, if any,
applied in the comparison (identity, rescaling, requantization), with its
own uncertainty contribution; (iv) the evidence establishing the
comparability---which, for anything beyond definitional identity, means a
measurement campaign under both profiles with agreement scored against
declared tolerances; (v) its own validity window and drift obligations,
since comparability decays for the same reasons calibration does; and (vi)
its own review state. In other words, a comparison rule is itself a
profile-like object: typed, evidenced, digested, expiring, reviewed. The
symmetry is intentional. If crossing between profiles were cheaper than
staying inside one, the system would leak at the crossings; making the
crossing artifact carry the same discipline as the endpoints closes the
arbitrage.

\begin{remark}\label{ch05:rem:comparison}
The prespecified comparison rule required by the registered
checkpoint-replication hypothesis is the intended first instance: a rule
relating a profile on one checkpoint to a profile on a held-out checkpoint,
admitted \emph{before} the replication runs, so that the replication cannot
quietly tune the comparison to the result. As of this writing no such rule
exists; the slot for it does.
\end{remark}

\section{Estimating Effective Resistance Under a Profile}\label{ch05:reff}

The profile exists so that one quantity can be estimated honestly: the
effective resistance $r_{\mathrm{eff}}$, in alpha-per-activation units, for
a given test point. The estimation requirements, stated here and developed
statistically in Chapter~\ref{ch06:response}, are:

\begin{enumerate}
\item \emph{Fit on the calibration split only.} The estimator consumes
dose-response observations drawn from the calibration split referenced by
the profile, at doses from the profile's alpha grid, under the profile's
seed and randomization method.
\item \emph{Declared method, receipted run.} The estimator is the profile's
named fit method; the run is recorded in the fit receipt with its
diagnostics.
\item \emph{Uncertainty mandatory.} The estimate carries an uncertainty
produced by the profile's named uncertainty method; a bare point estimate is
not an admissible resistance.
\item \emph{Region-conditioned.} The fit claims validity only in the active
region as classified by the profile's response-region method; observations
in cutoff, saturation, breakdown, or unstable regions do not certify it.
\item \emph{Predict held-out data.} The estimate's admissible role,
\texttt{CALIBRATED\_HELD\_OUT}, is earned by prediction on the validation
split---data the fit never saw---scored against the profile's tolerances.
Defining $r = a/y$ on a sample and ``verifying'' $\hat y = a/r$ on the same
sample is the circular procedure the candidate rejects outright.
\item \emph{Discrepancy acknowledged.} Following the computer-model
calibration literature, the estimation must represent not only parameter
uncertainty but model discrepancy---the structured difference between the
best-fitting linear response and the system's actual response---rather than
treating the best fit as truth \cite{kennedyohagan}. A calibration that
hides discrepancy inside its parameter estimate will predict confidently
and wrongly; one that models discrepancy knows where its linear hypothesis
stops.
\end{enumerate}

Every requirement lands on a profile field: (1) on the split references and
randomization fields; (2) on the fit method and receipt; (3) on the
uncertainty method; (4) on the response-region method; (5) on the distinct
validation split and the tolerance profile; (6) on the receipt's required
diagnostics.

Two interactions between estimation and the rest of the profile deserve
explicit note. First, estimation interacts with \emph{drift}: a resistance
fitted under a profile inherits the profile's validity window, and the
drift-check discipline is what stands between ``was calibrated'' and ``is
calibrated.'' A fitted $r_{\mathrm{eff}}$ under a profile whose drift state
has left \texttt{CURRENT} is not merely older evidence; it is evidence
whose preconditions are no longer certified, and the currency assertion
removes it from service automatically rather than trusting consumers to
check dates. Second, estimation interacts with \emph{region
classification} bidirectionally: the fit claims validity only in the
active region, but the active region is itself identified from sweep data
by the profile's region method. The candidate handles the circularity the
only defensible way---both procedures are frozen in the profile before
either runs, so neither can be adjusted to flatter the other, and the
region method's verdicts are replayable facts under the digest. This is the sense in which the profile is not bureaucracy
around the science but the science's own preconditions, itemized. The
chapter's boundary, however, is unchanged from the monograph's: these are
requirements an estimate \emph{would} have to meet. No estimate meeting
them---no admitted live calibration of any production system---exists in the
current evidence, and nothing in this chapter asserts otherwise.

\section{What the Profile Is Not}\label{ch05:isnot}

Because the profile is the most infrastructure-like object in the
candidate, it attracts misreadings that the chapter should close by naming.

\emph{The profile is not a certification.} Passing profile validation means
the record is well formed and self-consistent: references resolve, splits
differ, grids are ordered, coverage rules hold. It does not mean the
calibration is good, the direction meaningful, the corpus adequate, or the
model safe. The review state carries human judgment about sensibleness;
the authenticated-eligibility machinery (constant false in this candidate)
would carry independent authority; neither is implied by structural
validity, and nothing in the digest should be read as a seal of quality.

\emph{The profile is not a claim that its fields are sufficient.} The field
set was derived from the audit's conflicts and from general reproducibility
practice; it forecloses the enumerated failures of
Table~\ref{ch05:tab:threats}. Whether further conditions move numbers in
ways the profile does not record---hardware nondeterminism below the dtype
level, undocumented runtime kernel dispatch, environmental dependencies
outside the lock---is an open empirical question. The honest statement is
that the profile records every condition the candidate's authors currently
know to matter, and that discovering another such condition would mean
extending the schema, versioning it, and treating older profiles'
comparability to newer ones as---again---a comparison-rule question rather
than an assumption.

\emph{The profile is not portable authority.} A profile is bytes; anyone can
construct one. Its digest proves integrity (these exact fields, unaltered),
never provenance in the social sense (a trustworthy party constructed
this). The trust architecture that would bind profiles to accountable
parties---protected key lifecycles, independent authority manifests,
trusted time---is deliberately unimplemented, and the candidate's operation
contexts say so on every binding by carrying unauthenticated authorization
states. Treating a structurally valid profile from an unknown source as
trustworthy input would be a category error the type system permits but the
process rules forbid.

\emph{The profile is not the science.} It is the container the science
requires. The registered hypotheses of
Chapter~\ref{ch12:hypotheses} where present---local response above matched
controls, held-out predictive resistance, gain composition, redundancy,
residual localization---are the science, and every one of them presupposes a
profile precisely so that a failed prediction is informative: with all
conditions frozen and recorded, a miss cannot be blamed on an unrecorded
condition, which is what makes the hypotheses falsifiable in the first
place. A framework that let conditions float would be unfalsifiable in the
weakest sense---its failures would always have an alibi. The profile
removes the alibis, for successes and failures alike.

\begin{boundary}\label{ch05:bound}
This chapter specifies a data contract and its rationale. It establishes no
empirical adequacy of any field choice, no validated calibration of any live
system, no comparability of any pair of profiles, and no authority beyond
\textsc{ring 4 / not licensed / held}. Where the chapter argues that a field
prevents a failure, the claim is architectural---the failure becomes
unrepresentable or fail-closed in the candidate's types---not a claim that
deployed systems exhibiting the failure have been observed or measured.
\end{boundary}
